\documentclass[11pt]{article}
\usepackage{amsmath,amssymb,amsthm}
\usepackage[margin=1in]{geometry}
\newtheorem{proposition}{Proposition}
\title{P6: on the transcendence of $\beta_2$ (status note)}
\author{}
\date{2026-09-26}
\begin{document}
\maketitle

\noindent\textbf{Scope.} This note documents an attempt at research-programme item P6
(``$\beta_2$ transcendence''). It largely \emph{reproduces and cross-checks} an existing,
much more thorough internal investigation (\texttt{private/ROOM/Beta2/report\_G1.md} and
\texttt{private/ROOM/Beta2/g1\_certificate/}), rather than opening new ground: that room
already produced a machine-checked 2090-digit certified bracket and LLL/PSLQ exclusion
bounds. What follows restates the precise definitions from the actual paper/code, reproduces
the key numbers independently at lower precision, and gives an honest assessment of the
$\rho=1/2$ obstruction named in \texttt{private/RESEARCH\_PROGRAM\_OPEN.md} row P6.

Every claim below is labelled \textbf{PROVED}, \textbf{VERIFIED} (numerically/computationally,
by me in this session or reproducing a prior computation), or \textbf{HEURISTIC}. Anything
attributed to the literature (Tschakaloff, Bézivin--Bundschuh--Väänänen, André, Di~Vizio--Hardouin,
Koelink--Swarttouw, Mahler's method) is \emph{recalled from training, not re-checked against a
primary source in this session}, and is flagged as \textbf{UNVERIFIED-CITATION} accordingly.

\section{Precise definition of $\beta_2$}

From \texttt{paper/journal/paper2a.tex}, Corollary~\ref{cor:beta} (paper's own numbering,
reproduced here as prose):
\begin{quote}
The exponential growth rate of the sequences $u_n,v_n$ is
$\beta_2 = 1.4916177871\ldots$, where $\beta_2^{2}$ is the reciprocal of the smallest travel
pole $q_1\in(0,1)$; the growth-rate lower bound $\beta_2$ for $u_n,v_n$ is unconditional (free
even of hypothesis (M)); the value $3/2$ is excluded as the growth rate.
\end{quote}
and, separately (line 2672 of the same file):
\begin{quote}
Is the number $\beta_2$ transcendental? The question reduces to a value statement for a
function defined by \cite[Appendix app:beta2]{Bonfioli2026bB}, independent of everything else
in the paper.
\end{quote}

Concretely, $q^\ast=q_1\in(0,1)$ is the least positive root of the \emph{base-variable} power
series (paper2 rem:betanumber; reproduced verbatim in
\texttt{code/zeta\_probe/beta2\_effective\_irrationality.py} and
\texttt{private/ROOM/Beta2/g1\_certificate/step1\_bracket.py}):
\[
S(q) \;=\; 1 + \sum_{k\ge 1} (-1)^k \, q^{2k(k-1)} \, \frac{[2q(1-q)]^k}
{\prod_{j=1}^{k}(1-q^{2j-1})(1-q^{2j})}, \qquad q\in(0,1),
\]
and
\[
\beta_2 = 1/\sqrt{q^\ast}.
\]
Here ``base variable'' means $q$ itself is the free variable of the series (as in $q$-series /
basic hypergeometric theory), as opposed to a generating-function variable whose \emph{coefficients}
would be the arithmetic objects of interest; this is exactly the distinction the ledger's phrase
``base-variable zero'' is pointing at. \texttt{code/zeta\_probe/route\_b/beta2\_qniven\_ladder.py}
identifies $S(q)=G(q,2q(1-q))$ with a Hahn--Exton (3rd Jackson) $q$-Bessel function
$J^{(3)}_{-1/2}(\cdot;q^2)$ up to reparametrization, matching the $J^{(3)}_{-1/2}$ notation used in
paper2a's own statement of Corollary~\ref{cor:beta} \textbf{[VERIFIED against paper2a.tex and
the code, low precision, this session]}.

\textbf{What is already PROVED (unconditional):} $\beta_2$'s role as the exact exponential
growth rate of $u_n,v_n$, and the exclusion of $3/2$ as that rate (memory note
\texttt{route-b-singularity}, cor:beta). This is a statement about a growth \emph{rate}, not
about the arithmetic nature of the real number $\beta_2$.

\textbf{What is OPEN:} whether the number $\beta_2$ (equivalently $q^\ast$) is irrational, or
transcendental. This is the P6 target.

\section{Numerics: independent reproduction}

\textbf{VERIFIED (this session, mpmath, 320 decimal digits, Newton's method on $S$, wrapped in
the required \texttt{perl alarm 290} cap, script
\texttt{code/zeta\_probe/rootunity\_zeros/general/P6/p6\_verify.py}):}
\[
q^\ast = 0.449453630558948046125545825395706089225112808545993877680333\ldots
\]
\[
\beta_2 = 1.49161778711437422683671274698306319965329579126010371956739\ldots
\]
These match, digit for digit, the values already on record in
\texttt{private/ROOM/Beta2/g1\_certificate/bracket.txt} (a pre-existing 2090-digit certified
bracket) and in memory note \texttt{route-b-singularity}. Residual $|S(q^\ast)|\approx
5\times10^{-322}$ at 320-digit working precision, consistent with convergence to the true root
rather than an artifact.

\textbf{VERIFIED (this session, mpmath PSLQ, 320-digit precision, small fixed basis):} no
integer relation was found by PSLQ (maxcoeff $10^6$, maxsteps 2000) among
$\{1,\ \beta_2,\ \beta_2^2,\ \beta_2^3,\ \pi,\ e,\ \log 2,\ \sqrt2,\ \varphi\}$ ($\varphi$ = golden
ratio), nor (maxcoeff $10^8$, maxsteps 3000) among $\{1,\beta_2,\ldots,\beta_2^6\}$ alone. This
is a much weaker check than the existing certificate below, included only as an independent
cross-check with a different tool run in this session.

\textbf{Pre-existing, stronger result (VERIFIED, reproduced from
\texttt{private/ROOM/Beta2/g1\_certificate/}, adversarially reviewed per
\texttt{report\_G1.md}, not re-run in full in this session for compute-budget reasons ---
LLL-reduction to degree 128 is well outside the 290s/3GB caps in force here):} using a
rigorous LLL certificate (exact Gram--Schmidt norms on the LLL-reduced basis via
\texttt{python-flint}, from a 2090-digit certified interval bracket for $q^\ast$ and for
$\beta_2$),
\begin{itemize}
\item $q^\ast$ is not a root of any integer polynomial of degree $\le 16$ with height
$<10^{120}$, nor degree $\le 8$ with height $<10^{230}$, nor degree $\le 64$ with height
$<10^{28}$;
\item $\beta_2$ obeys comparable bounds;
\item a planted negative control ($\sqrt2-1$, a genuine degree-2 algebraic number) is
correctly \emph{not} excluded by the same method, i.e.\ the method does not over-claim.
\end{itemize}
These are far stronger than anything re-derivable in this session's compute budget, and I have
not found any flaw in \texttt{step1\_bracket.py} / \texttt{step2\_lattice.py} on inspection; I
have not, however, independently re-executed the degree-64/128 LLL runs myself this session, so
I report them as \textbf{VERIFIED-BY-PRIOR-SESSION}, not re-verified here.

\section{The $\rho=1/2$ obstruction: removable or fundamental?}

\texttt{code/zeta\_probe/route\_b/beta2\_qniven\_ladder.py} defines $\rho$ as the ratio of
(archimedean decay rate supplied by the numerator, here $q^{k^2}$-type terms) to (growth rate of
the cleared denominators, here the double $q$-Pochhammer $(q;q)_{2k}$) in the relevant
$q$-continued-fraction / Padé / Lommel-ladder construction for $S$ (equivalently for the
Hahn--Exton $q$-Bessel function identified with it). Concretely, several independent
constructions were tried and each lands at the same ratio:
\begin{itemize}
\item $z$-Padé approximants to $S$: remainder $\sim q^{3.5n^2}$ against denominator $q$-degree
$\sim 9n^2$;
\item the $z$-continued fraction (Pincherle limit = ratio of the two Hahn--Exton solutions):
geometric numerator decay vs.\ $t^{n^2}$ denominator growth;
\item the $\nu$-ladder (q-Lommel/Niven recursion) at the zero $x_0$: $J_{m+1/2}(x_0)\sim
Cx_0^m$ is only \emph{geometric} in $m$ (because the $q$-factorial that would normally supply
super-geometric decay is itself only geometric, unlike the classical integer-factorial Niven
setting), against denominators growing like $t^{m^2}$;
\item the base$\leftrightarrow$argument ($b\leftrightarrow z$) swap identity for the
${}_1\phi_1$ preserves $\rho$ rather than improving it;
\item a Hankel-determinant / resultant construction on the moment sequence of $\beta_2$
produces denominators growing \emph{cubically} in the exponent, which is worse, not better
(\texttt{beta2\_hankel\_nogo.py}: ``deepens the $\rho=1/2$ deficit'').
\end{itemize}
All of the above are \textbf{HEURISTIC/EXPLORATORY} results from earlier sessions (I did not
re-derive the Padé remainder estimates or the Lommel-ladder degree count myself this session);
I read the code and find the stated degree/order bookkeeping internally consistent, but flag
that I have not independently re-verified e.g.\ the claim ``$\deg_q p_m = m^2$ exactly'' by
symbolic computation in this session.

\textbf{Is $\rho=1/2$ a coincidence of the particular series representation, or forced by the
object itself?} My assessment, based on reading five independent constructions (Padé, continued
fraction, Lommel ladder, base/argument swap, Hankel) all landing on the \emph{same} ratio: this
looks \textbf{structural, not an artifact of one bad parametrization}. The double Pochhammer
$(q;q)_{2k}$ in the denominator of $S$'s own defining series already fixes the growth rate of
any naturally associated integer-coefficient approximation scheme, and every reparametrization
tried preserves or worsens that rate rather than improving it. This is consistent with (and, I
believe, is really the same fact as) the general phenomenon that $q$-Bessel/$q$-exponential
functions of this Hahn--Exton type sit exactly at the boundary of applicability of the classical
Tschakaloff/Bézivin--Bundschuh--Väänänen $q$-irrationality machinery, which (as I recall it, and
flag as \textbf{UNVERIFIED-CITATION}, not re-checked against a primary source this session)
needs a decay-vs-denominator-growth ratio $\rho\ge1$ to run a Siegel-type linear-independence
argument. I could not find, and did not construct, any transformation of $S$ that changes the
$1/2$ to something $\ge1$; I regard this as evidence toward ``fundamental'' rather than
``removable,'' but it falls short of a proof that no such transformation exists.

\textbf{Base-variable vs.\ coefficient-variable theories.} Separately, \texttt{report\_G1.md}'s
verdict --- ``there is no base-variable theory for $q$-series that are neither modular nor
Mahler'' --- points at a second, independent gap: $\beta_2$'s defining object $S(q)$ is a
statement about the value of a function \emph{at} a specific base point $q^\ast$ (a
transcendence-of-values-of-a-$q$-series-in-its-own-base question), not a statement about the
arithmetic nature of Taylor coefficients of a fixed function (which is what Mahler's method and
most modular-function transcendence results address). \texttt{beta2\_diffalg\_search.py} and
\texttt{beta2\_ade\_sweep.py} (not re-run this session; large prior body of exploratory work)
apparently searched for a differential-algebraic or modular reformulation and did not find one
that would let a standard theory apply; I did not re-verify these negative results myself.

\section{Strongest honest statement}

Combining what I re-verified this session with what the prior room verified and I did not
independently re-run:

\begin{proposition}[status, mixed labels]
\begin{enumerate}
\item[(a)] \textbf{PROVED} (paper2a, Corollary cor:beta): $\beta_2=1.4916177871\ldots$ is the
exact, unconditional exponential growth rate of $u_n,v_n$; $3/2$ is excluded as that rate.
\item[(b)] \textbf{VERIFIED} (this session, 320 digits, and independently in the prior room to
2090 digits): $q^\ast=\beta_2^{-2}=0.44945363055894804612554582539\ldots$ is not a root of any
integer polynomial of degree $\le16$ and height $<10^{120}$ (nor several other degree/height
combinations listed in \S2); the same holds for $\beta_2$ itself.
\item[(c)] \textbf{HEURISTIC}: the obstruction to promoting (b) into a proof of irrationality
(let alone transcendence) is a fixed ratio $\rho=1/2$ between archimedean decay and denominator
growth across every natural approximation scheme tried for the Hahn--Exton $q$-Bessel value
$S(q^\ast)=0$; existing $q$-irrationality theorems need $\rho\ge1$. This looks structural rather
than removable, but that is an assessment, not a proof of non-removability.
\item[(d)] \textbf{OPEN}: irrationality of $\beta_2$; transcendence of $\beta_2$.
\end{enumerate}
\end{proposition}

\textbf{Correction (post-review, Reviewer AJ):} \texttt{report\_G1.md} itself distinguishes two
tiers in its LLL exclusion runs: degree $\le16$ (height $<10^{120}$), $\le8$ (height $<10^{230}$)
and $\le64$ (height $<10^{28}$) are \textbf{PROVED}, backed by an explicit unimodularity check
(\texttt{assert T*B==L and abs(det T)==1} in \texttt{step2\_lattice.py}); the degree-24, -96 and
-128 runs predate that check and are labeled \textbf{VERIFIED only} there, a strictly weaker
guarantee. The sentence below is corrected to scope the PROVED-tier claim to degree $\le64$ and
report degree 128 separately, at its true (weaker) tier.

The strongest \emph{true} statement I can make, in one sentence: \emph{$\beta_2$ is a
well-defined, unconditionally-established growth-rate constant that is provably (PROVED tier) not
a root of any integer polynomial of degree $\le64$ under the height caps in \S2, and additionally
(VERIFIED-only tier, predating the prior room's unimodularity check) not one of degree $\le128$
under a further height cap, but no proof of even its irrationality is known, and the one natural
route (a $q$-Siegel/Shidlovskii-type argument on its Hahn--Exton $q$-Bessel representation) is
blocked at exactly $\rho=1/2$, a factor of two below every published sufficient condition I am
aware of.}

\section{What would be needed for full transcendence}

\begin{itemize}
\item A $q$-irrationality (or $q$-Siegel--Shidlovskii-type transcendence) criterion valid at
$\rho=1/2$ rather than $\rho\ge1$ --- i.e., a genuinely new arithmetic input, not a
reparametrization of $S$ (five reparametrizations were tried and none moved $\rho$).
\item Alternatively, a value-transcendence theory for $q$-series transcendental \emph{in the
base variable} that are neither modular functions nor Mahler functions (report\_G1.md's
diagnosis of the missing theory).
\item Independently: a proof (not merely absence of a modular/Mahler/D-finite structure found
by search) that $S(q)$, viewed as a function of $q$, is not annihilated by any algebraic
differential or $q$-difference equation of a type where such a theory does exist; the searches
in \texttt{beta2\_diffalg\_search.py}/\texttt{beta2\_ade\_sweep.py} are negative but exploratory,
not proofs of non-existence.
\end{itemize}

\section{Weak points of this note}

\begin{itemize}
\item The literature citations (Tschakaloff, Bézivin--Bundschuh--Väänänen, André,
Di~Vizio--Hardouin, Koelink--Swarttouw's exact identity numbers) are recalled from training
data and were \emph{not} checked against primary sources in this session; they may be
mis-stated or mis-numbered.
\item The degree-64/128 LLL certificates and the Padé/Lommel-ladder degree bookkeeping
(\S3) are taken from the prior room's outputs and code inspection, not re-executed here; I
verified internal consistency by reading, not by rerunning at full compute.
\item My own PSLQ check (\S2) is numerically much weaker than the existing certificate and
is included only as an independent second opinion with a different tool, not as new evidence
beyond what already existed.
\item The claim that $\rho=1/2$ is ``structural'' (\S3) is an inductive judgment from five
failed reparametrizations, not a proof that no sixth reparametrization could do better.
\end{itemize}

\end{document}
